\subsection{正交基。}

定义 1. --- 设 Φ 是 E 上的一个埃尔米特型。若 E 的一个基 ( \(e_{i}\) ) 中任何两个元素都对 Φ 正交，则称这个基对 Φ 是正交的。

此外，若有 \(\Phi(e_i, e_i) = 1\) （对一切 \(i\) ），则称基 \((e_i)\) 是规范正交的。

设 \((e_i)\) 是一个正交基；若令 \(\Phi(e_i, e_i) = \alpha_i\) ，则有

\[
\Phi (\sum_ {i} \xi_ {i} e _ {i}, \sum_ {i} \eta_ {i} e _ {i}) = \sum_ {i} \xi_ {i} \alpha_ {i} \overline {{\eta}} _ {i}.\tag{1}
\]

引理 1. --- 假设 A 是一个体，\(\Phi\) 是 E 上的一个埃尔米特型 \(\neq 0\) 。若 E 的所有向量都是迷向的，则 A 是特征 2 的交换域，反自同构 J 是恒等映射，且 \(\Phi\) 是交错的。

实际上，把 \(\Phi(x+y, x+y)=0\) 展开，并考虑到假设 \(\Phi(x,x)=\Phi(y,y)=0\) ，就得到关系 \(\Phi(x,y)=-\overline{\Phi(x,y)}\) （对 E 中所有 \(x,y\) ）。由于 \(\Phi\) \(\neq 0\) ，E 中存在 \(x,y\) 使 \(\Phi(x,y)=1\) 。写出 \(\Phi(\lambda x,y)=-\overline{\Phi(\lambda x,y)}\) ，就得到 \(\bar{\lambda}=-\lambda\) （对一切 \(\lambda\in A\) ）。先取 \(\lambda=1\) ，就看出 A 的特征是 2；于是关系 \(\bar{\lambda}=-\lambda\) 表明 J 是恒等映射，从而 A 是交换的且 \(\Phi\) 是交错的。

定理 1. --- 假设 A 是一个体，E 是 A 上有限维数为 n 的向量空间。则对 E 上的每个埃尔米特型 Φ，E 都容许一个正交基，除非下列条件同时满足：

\begin{enumerate}
\def\labelenumi{(\Alph{enumi})}
\setcounter{enumi}{2}
\tightlist
\item
  A 是特征 2 的交换域，反自同构 J 是恒等映射，Φ 是交错的且非零。
\end{enumerate}

对 \(n\) 作归纳推理，因为 \(n=0\) 时结果是显然的。我们可以假设 \(\Phi\neq0\) 。若 (C) 不满足，引理 1 表明存在一个元素 \(x\in\mathbf{E}\) 使 \(\Phi(x,x)\neq0\) 。设 H 是 E 中与 \(x\) 正交的子空间；它的维数 \(\geqslant n-1\) ，并且由于 \(x \notin H\) ，\(H\) 的维数恰好是 \(n-1\) 。若限制 \(\Psi\) （把 \(\Phi\) 限制到 \(H\) 上所得）不满足 (C)，则依归纳假设，存在一个基 \((e_2, \ldots, e_n)\) （\(H\) 的），它对 \(\Psi\) 是正交的；令 \(e_1 = x\) ，就得到一个正交基 \((e_1, e_2, \ldots, e_n)\) （\(E\) 的）。剩下要考察的是下述情形：\(A\) 是特征 2 的交换域，\(J\) 是恒等映射，且 \(\Psi\) 交错而 \(\neq 0\) 。此时存在 \(y\) 、\(z\) （\(H\) 中的）使 \(\Psi(y, z) \neq 0\) ；令 \(e_1 = x + y\) ；要使 \(x + \lambda z (\lambda \in A)\) 与 \(e_1\) 正交，必须且只需有 \(0 = \Phi(x + y, x + \lambda z) = \Phi(x, x) + \lambda \Psi(y, z)\) ，这个条件以唯一的方式确定 \(\lambda\) ；这样选定数量 \(\lambda\) 后，我们有 \(\Phi(x + \lambda z, x + \lambda z) = \Phi(x, x) \neq 0\) ，因而限制 \(\Psi'\) （把 \(\Phi\) 限制到 \(H'\) （\(E\) 的与 \(e_1\) 正交的子空间）上所得）不是交错的；于是可以对 \(H'\) 应用归纳假设，这就证明了定理。

当 (C) 满足时，显然不存在对 \(\Phi\) 的正交基。

推论 1. --- 沿用定理 1 的记号，再假设 (C) 不满足，\(\Phi\) 非退化，并且对每个 \(x \in E\) ，存在 \(\rho \in A\) 使 \(\Phi(x, x) = \rho \bar{\rho}\) 。则 E 容许一个对 \(\Phi\) 的规范正交基。

实际上，设 \((e_i)\) ( \(i = 1, \ldots, n\) ) 是 E 的一个正交基。令 \(\Phi(e_i, e_i) = \alpha_i\) 。我们有 \(\alpha_i \neq 0\) （对 \(i = 1, \ldots, n\) ），因为 \(\Phi\) 非退化。根据假设，存在 A 的元素 \(\beta_i\) ，使得 \(\alpha_i = \beta_i \overline{\beta_i}\) （对 \(i = 1, \ldots, n\) ）；我们有 \(\beta_i \neq 0\) 。令 \(f_i = \beta_i^{-1} e_i\) ，则有 \(\Phi(f_i, f_i) = \beta_i^{-1} \alpha_i \overline{\beta_i^{-1}} = 1\) （对一切 \(i\) ），且 \(\Phi(f_i, f_j) = 0\) （对 \(i \neq j\) ）。于是 \((f_i)\) 是一个规范正交基。

注. --- 当 J 是恒等映射且 A 的每个元素都是 A 中某个元素的平方时（例如当 A 代数闭时），推论的最后假设成立。

推论 2. --- 设 A 是一个体，R 是 A 上一个 n 阶秩 r 的埃尔米特矩阵。那么，除非下列条件成立：

(C') A 是特征 2 的交换域，J 是恒等映射，R 是交错的且非零，

存在 A 上的一个 n 阶可逆矩阵 P，使得

\[
{ } ^ { t } P . R . \overline { { P } } = \left( \begin{array} { c c c c c } \alpha _ { 1 } & 0 \dots 0 \dots 0 \\ 0 & \alpha _ { 2 } \dots 0 \dots 0 \\ \cdot & \cdot & \cdot & \cdot & \cdot \\ 0 & 0 \dots \alpha _ { r } \dots 0 \\ 0 & 0 \dots 0 \dots 0 \\ \cdot & \cdot & \cdot & \cdot & \cdot \\ 0 & 0 \dots 0 \dots 0 \end{array} \right)
\]

\(where \overline{\alpha}_{i} = \alpha_{i} \neq 0 \text{ 对于 } i = 1, \ldots, r.\)

命题 1. --- 假设 A 是一个交换域。设 \(\Phi\) 是 E 上的一个埃尔米特型，\((x_{n})\) ( \(n=1,2,\ldots\) ) 是 E 的线性无关向量的一个序列（有限或无限），使得对每个 \(n\) ，子空间 \(\mathrm{E}_{n}=\mathrm{A}x_{1}+\cdots+\mathrm{A}x_{n}\) 都是非迷向的。以 \(\mathrm{D}_{jn}\) ( \(j\leqslant n\) ) 表示 \(\Phi(x_{j},x_{n})\) 在矩阵 \((\Phi(x_{s},x_{t}))_{(s,t=1,\ldots,n)}\) 中的代数余子式。于是我们有 \(\mathrm{D}_{nn}\neq0\) （对一切 \(n\) ）。令

\[
e _ {n} = \sum_ {j = 1} ^ {n} \mathrm{D} _ {n n} ^ {- 1} \mathrm{D} _ {j n} x _ {j}.\tag{2}
\]

则对每个 \(n\) ，\((e_{1},\ldots,e_{n})\) 是 \(\mathrm{E}_{n}\) 的一个正交基，并且我们有

\[
\Phi (e _ {n}, e _ {n}) = \mathrm{D} _ {n n} ^ {- 1} \mathrm{D} _ {n + 1, n + 1}.\tag{3}
\]

实际上，由于 \(\Phi\) 在 \(\mathbf{E}_{n-1}\) 上的限制非退化，我们有 \(\mathrm{D}_{nn} \neq 0\) ( \(\S 2\) ，命题 3)；注意我们有 \(\mathrm{D}_{11} = 1\) ，因为空矩阵的行列式等于 1。公式 (2) 首先蕴含 \(e_n \equiv x_n\) (mod. \(\mathbf{E}_{n-1}\) ) （对一切 \(n\) ），从而诸 \(e_n\) 线性无关，且 \((e_1, \ldots, e_n)\) 是 \(\mathbf{E}_n\) 的一个基。对每个 \(j < n\) ，有

\[
\Phi (e _ {n}, x _ {j}) = \mathrm{D} _ {n n} ^ {- 1} \sum_ {k = 1} ^ {n} \mathrm{D} _ {k n} \Phi (x _ {k}, x _ {j}) = 0
\]

（第 III 章，§ 6，第 1 小节，公式 (12)）；因而 \(e_n\) 与 \(\mathrm{E}_{n-1}\) 正交，特别地，与 \(e_j\) （对 \(j < n\) ）正交。另一方面，我们有

\[
\begin{array}{r l} \Phi (e _ {n}, e _ {n}) & = \Phi (e _ {n}, \sum_ {j = 1} ^ {n} \mathrm{D} _ {n n} ^ {- 1} \mathrm{D} _ {j n} x _ {j}) = \Phi (e _ {n}, x _ {n}) = \Phi (\sum_ {j = 1} ^ {n} \mathrm{D} _ {n n} ^ {- 1} \mathrm{D} _ {j n} x _ {j}, x _ {n}) \\ & = \mathrm{D} _ {n n} ^ {- 1} \sum_ {j = 1} ^ {n} \mathrm{D} _ {j n} \Phi (x _ {j}, x _ {n}) = \mathrm{D} _ {n n} ^ {- 1} \mathrm{D} _ {n + 1, n + 1} \end{array}
\]

（第 III 章，§ 6，第 1 小节，公式 (10)）。这就证明了我们的断言。

利用命题 1 的记号，我们称序列 \((e_n)\) 是由序列 \((x_n)\) 通过格拉姆–施密特正交化程序得到的。

命题 2. --- 设 \(\Phi\) 是 E 上的一个埃尔米特型，\((e_i)\) ( \(i = 1, \ldots, n\) ) 是 E 的一个对 \(\Phi\) 的正交基（相应地，规范正交基）。则对每个 \(p \geqslant 0\) ，\(\bigotimes^{p} \mathbf{E}\) 的由诸 \(e_{i_1} \otimes \cdots \otimes e_{i_p}\) 形成的基以及基 \((e_h)\) （\(\bigwedge^{p} \mathbf{E}\) 的）（其中 H 跑遍具有 \(p\) 个元素的子集（取自 {[}1, n{]} ）所成的集；参见第 III 章，§ 5，第 6 小节），对 \(\Phi\) 到 \(\bigotimes^{p} \mathbf{E}\) 与 \(\bigwedge^{p} \mathbf{E}\) 上的扩张（分别地）是正交的（相应地，规范正交的）（§ 1，第 9 小节）。此外，若与 \(\Phi\) 相伴的映射都是双射的，则基 \((e_i')\) （\(\mathbf{E}^*\) 的、与 \((e_i)\) 对偶的）对逆形式 \(\hat{\Phi}\) （\(\Phi\) 的）（§ 1，第 7 小节）是正交的（相应地，规范正交的）。

关于 \(\widehat{\otimes}\) E 与 \(\wedge\) E 的断言由 § 1 第 9 小节的公式 (35) 与 (37) 立即得出。关于逆形式的断言则由下述事实得出：\(\widehat{\Phi}\) 关于 \((e_i')\) 的矩阵是 \(\Phi\) 关于 \((e_i)\) 的矩阵的逆矩阵 ( \(§ 1\) ，第 10 小节)。

